\documentclass[12pt]{article}
\usepackage[T1]{fontenc}
\usepackage[utf8]{inputenc}
\usepackage{lmodern}
\usepackage{textcomp}
\usepackage{enumitem}
\usepackage{geometry}
\geometry{margin=1in}
\begin{document}
\begin{center}
Answer Key - Economics - Undergraduate Level Assessment 22 \
\small (Answers are ordered exactly as in the question paper.)
\end{center}

\section*{Multiple Choice}
\begin{enumerate}[label=\arabic*.]
\item \textbf{Answer:} A
\\ \textit{Options:} A) the real quantity of money increases; B) imports become relatively more expensive; C) the value of cash increases; D) exports become relatively cheaper; E) inflation rate decreases
\\ \textit{Note:} The correct answer is based on the liquidity effect, which states that an increase in the price level reduces the real value of money balances, leading to a decrease in consumer spending and investment, thereby resulting in a downward-sloping aggregate demand curve.
\item \textbf{Answer:} A
\\ \textit{Options:} A) 1 and 0.5; B) 2 and 0.5; C) -1 and -0.5; D) -1 and 2; E) 3 and -2
\\ \textit{Note:} The characteristic roots of the MA process are derived from the characteristic equation formed by the coefficients of the lagged error terms, which in this case leads to the roots being 1 and 0.5, indicating the influence of past shocks on the current value of the series.
\item \textbf{Answer:} A
\\ \textit{Options:} A) Chicago school, Stiglitz; B) Marxist school, Minsky; C) Behavioral school, Akerlof; D) Neoclassical school, Friedman; E) Neoclassical school, Coase
\\ \textit{Note:} The Chicago school of economists supports Say's Law, which posits that supply creates its own demand, while Joseph Stiglitz is a prominent economist who has criticized this theory, arguing that demand deficiencies can lead to prolonged economic downturns.
\item \textbf{Answer:} D
\\ \textit{Options:} A) a decrease in the money supply.; B) an increase in consumer confidence.; C) a sudden surge in technological innovation.; D) a prolonged period of very bad weather.; E) a decrease in consumer confidence.
\\ \textit{Note:} A prolonged period of very bad weather can disrupt agricultural production and supply chains, leading to higher prices (inflation) and reduced output (lower real GDP), which aligns with the scenario of increased price levels and decreased real GDP in the short run.
\item \textbf{Answer:} A
\\ \textit{Options:} A) The current value of y; B) The median value of y over the in-sample period; C) The minimum value of y over the in-sample period; D) The value of y two periods back; E) One
\\ \textit{Note:} In a random walk, the best predictor of the next value is the current value because the future path of the series is unpredictable and solely dependent on its most recent observation.
\end{enumerate}

\section*{True / False}
\begin{enumerate}[label=\arabic*.]
\setcounter{enumi}{5}
\item \textbf{Answer:} True
\\ \textit{Options:} A) True; B) False
\\ \textit{Note:} The statement is true because with a marginal propensity to consume of 0.85, consumers will spend 85\% of the $10 billion tax cut (which is $8.5 billion), leaving 15\% (or \$1.5 billion) to be saved.
\item \textbf{Answer:} True
\\ \textit{Options:} A) True; B) False
\\ \textit{Note:} The levels of incomes in foreign nations significantly influence the volume of a nation's exports and imports because higher income levels typically lead to increased demand for goods and services, affecting trade patterns between countries.
\item \textbf{Answer:} True
\\ \textit{Options:} A) True; B) False
\\ \textit{Note:} Cabbages grown in a summer garden are considered non-market production and are not included in GDP because GDP measures the value of goods and services that are sold in the marketplace.
\item \textbf{Answer:} False
\\ \textit{Options:} A) True; B) False
\\ \textit{Note:} Sales promotions in retail indicate a competitive market environment where firms use discounts and special offers to attract customers and differentiate themselves, rather than adhering to strict price collusion.
\item \textbf{Answer:} False
\\ \textit{Options:} A) True; B) False
\\ \textit{Note:} The statement is false because a leftward shift in the aggregate demand curve typically results in higher unemployment and lower inflation, aligning with the Phillips curve's implication of an inverse relationship between inflation and unemployment.
\end{enumerate}

\section*{Long Form Response}
\begin{enumerate}[label=\arabic*.]
\setcounter{enumi}{10}
\item \textbf{Answer:} To optimize resource allocation, the firm should decrease its capital and increase its labor. This adjustment is based on the relationship between the prices of labor and capital and their respective marginal products. By reducing capital, the firm can lower the marginal product of capital, making it less costly relative to labor. Simultaneously, increasing labor will raise its marginal product, enhancing productivity per unit of labor. This strategy allows the firm to achieve a more efficient allocation of resources, ultimately maximizing output while minimizing costs.
\\ \textit{Note:} correct because it emphasizes the need for the firm to realign its resource allocation by decreasing capital and increasing labor in response to their respective marginal products and prices, which leads to greater productivity and cost efficiency.
\item \textbf{Answer:} The aggregate supply curve can shift due to several factors, including changes in production costs, technological advancements, government policies, and the availability of resources. For instance, a decrease in the price of raw materials can shift the curve to the right, indicating an increase in supply at every price level. Conversely, an increase in the price level itself does not shift the aggregate supply curve; instead, it leads to a movement along the curve. This is because the price level reflects the prices of goods and services without altering the fundamental production capabilities or costs in the economy. Thus, while the price level may change the quantity supplied, it does not affect the overall supply capacity of the economy.
\\ \textit{Note:} correct because it accurately identifies that shifts in the aggregate supply curve are influenced by factors like production costs and technology, while changes in the price level result in movements along the curve rather than shifts.
\end{enumerate}

\vfill
\noindent\textit{End of Answer Key}
\end{document}